\section{Solving for \texorpdfstring{$\Lambda$}{Lambda} Algebraically}
\label{apd:solving-v}
Here we present the detailed derivation of \refeqn{cinvk}, including the linear system that is used to calculate $\Lambda$ from $v_*$, $C$, and $k_*$.  This derivation is included for clarity and completeness and is a review of the classical solution of least-squares with equality constraints as applied to our setting, together with the rank-one update rule that was proposed in \citet{bau-rewriting}.

We assume that $W$ is the optimal least-squares solution for memorizing a mapping from a previous set of keys $K$ to values $V$; this solution can be written using the normal equations as follows.
\begin{align}
\text{the $W$ that minimizes} \quad &|| W K - V ||_F^2 \\
\text{solves} \quad    & W K K^T = V K^T
\lbleq{normal-eq}
\end{align}
Here the Frobenius norm is used to write the total square error since the variable being optimized is a matrix $W$ rather than a vector $x$ as in the classical textbook presentation of least squares.

We wish to find a new matrix $\hat{W}$ that solves the same least squares problem with an additional equality constraint as written in \refeqn{cinvk}:
\begin{align}
    \hat{W}k_* = v_*
    \lbleq{eq-constraint}
\end{align}

This is the well-studied problem of least squares with a linear equality constraint.  The direct solution can be derived by defining and minimizing a Lagrangian, where $\Lambda \in \mathbb{R}^{H}$ minimizes the following:
\begin{align}
\text{define}\quad L(\hat{W}, \Lambda) &= \frac{1}{2} ||\hat{W}K-V||_F^2 - \Lambda^T(\hat{W}k_*-v_*) \\
&=\frac{1}{2} (\hat{W}K)(\hat{W}K)^T - V(\hat{W}K)^T + \frac{1}{2}VV^T - \Lambda^T(\hat{W}k_*-v_*) \\
\text{setting} \quad 0 = \frac{\partial L}{\partial \hat{W}} &= \hat{W}(KK^T) - VK^T -  \Lambda k_*^T\\
\hat{W}KK^T &= VK^T + \Lambda k_*^T
\lbleq{cls-normal-eq}
\end{align}
Subtracting \refeqn{normal-eq} from \refeqn{cls-normal-eq}, most terms cancel, and we obtain the update rule:
\begin{align}
    (\hat{W}-W) KK^T & = \Lambda k_*^T \\
    \hat{W} & = W + \Lambda (C^{-1}k_*)^T
    \lbleq{cls-soln}
\end{align}
The last step is obtained by defining $C = KK^T$, assuming $C$ is nondegenerate, and exploiting the symmetry of $C$.  Here we also write the row vector term as $u^T = (C^{-1}k_*)^T \in \mathbb{R}^{D}$, so we can write simply (rearranging \refeqn{cinvk} and \refeqn{cls-soln}):
\begin{align}
\hat{W}I - \Lambda u^T &= W
\lbleq{rome2}
\end{align}
To solve for $\Lambda$, we note that \refeqn{rome2} and \refeqn{eq-constraint}  form a linear system that allows both $\hat{W}$ and $\Lambda$ to be solved simultaneously if written together in block form.
\begin{align}
\left[\begin{array}{@{}c|c@{}}
\\
\quad \hat{W} \quad\, & \Lambda \\
\,
\end{array}\right]
\left[\begin{array}{@{}c|c@{}}
\\
\quad\; I \quad\;\; & k_* \\
\\ \hline
-u^T \rule{0pt}{2.2ex} & 0
\end{array}\right]
& =
\left[\begin{array}{@{}c|c@{}}
\mathstrut \\
\quad W \quad\, & v_* \\
\mathstrut 
\end{array}\right]
\end{align}
That is equivalent to substituting \refeqn{cls-soln} into \refeqn{eq-constraint} and calculating the following:
\begin{align}
    \hat{W} k_*
    & = (W + \Lambda u^T) k_*
    = W k_* + \Lambda (u^T k_*) = v_* \\
    \Lambda & = \frac{v_* - W k_*}{u^T k_*}
    = \frac{v_* - W k_*}{(C^{-1} k_*)^T k_*}
\lbleq{block-solution}
\end{align}
\clearpage